\section{Gradient, divergence, curl, and integral theorems}

This section develops the differential operator $\nabla$ (del) and its geometric meaning. In Cartesian coordinates,
\begin{equation}
\mathrm{d}f=\frac{\partial f}{\partial x}\,\mathrm{d}x
+\frac{\partial f}{\partial y}\,\mathrm{d}y
+\frac{\partial f}{\partial z}\,\mathrm{d}z.
\end{equation}

\subsection{Gradient}
For a scalar field $f(x,y,z)$, the gradient is
\begin{equation}
\nabla f=\frac{\partial f}{\partial x}\mathbf{e}_x
+\frac{\partial f}{\partial y}\mathbf{e}_y
+\frac{\partial f}{\partial z}\mathbf{e}_z.
\end{equation}
If
\begin{equation}
\mathrm{d}\boldsymbol{\ell}=\mathrm{d}x\,\mathbf{e}_x
+\mathrm{d}y\,\mathbf{e}_y
+\mathrm{d}z\,\mathbf{e}_z,
\end{equation}
then
\begin{equation}
\mathrm{d}f=\nabla f\cdot\mathrm{d}\boldsymbol{\ell}.
\end{equation}
The symbol $\nabla$ has no standalone meaning; combined with a scalar field, a dot product, or a cross product, it produces the gradient, divergence, or curl respectively.

Since
\begin{equation}
\mathrm{d}f=|\nabla f|\,|\mathrm{d}\boldsymbol{\ell}|\cos\theta,
\end{equation}
$\nabla f$ points in the direction of steepest increase of $f$. It is perpendicular to an equipotential surface, because motion along such a surface gives $\mathrm{d}f=0$.

In cylindrical coordinates $(r,\phi,z)$ and spherical coordinates $(r,\theta,\phi)$, the gradient becomes
\begin{equation}
\nabla f=\mathbf{e}_r\frac{\partial f}{\partial r}
+\mathbf{e}_\phi\frac{1}{r}\frac{\partial f}{\partial\phi}
+\mathbf{e}_z\frac{\partial f}{\partial z},
\end{equation}
and
\begin{equation}
\nabla f=\mathbf{e}_r\frac{\partial f}{\partial r}
+\mathbf{e}_\theta\frac{1}{r}\frac{\partial f}{\partial\theta}
+\mathbf{e}_\phi\frac{1}{r\sin\theta}\frac{\partial f}{\partial\phi}.
\end{equation}

\begin{example}[Gradients in different coordinates]
Let $a$ and $b$ be constants.

\textbf{(a) Cartesian coordinates.} For $f=ax^2y+by^2z$,
\begin{equation}
\nabla f=2axy\,\mathbf{e}_x+(ax^2+2byz)\,\mathbf{e}_y+by^2\,\mathbf{e}_z.
\end{equation}

\textbf{(b) Cylindrical coordinates.} For
$f=ar^2\sin\phi+ brz\cos 2\phi$,
\begin{equation}
\begin{aligned}
\nabla f
&=(2ar\sin\phi+bz\cos2\phi)\,\mathbf{e}_r\\
&\quad +(ar\cos\phi-2bz\sin2\phi)\,\mathbf{e}_\phi
+br\cos2\phi\,\mathbf{e}_z.
\end{aligned}
\end{equation}

\textbf{(c) Spherical coordinates.} For
$f=a+br\sin\theta\cos\phi$,
\begin{equation}
\nabla f=b\sin\theta\cos\phi\,\mathbf{e}_r
+b\cos\theta\cos\phi\,\mathbf{e}_\theta
-b\sin\phi\,\mathbf{e}_\phi.
\end{equation}
\end{example}

\subsection{Line integrals and path independence}
The work done by a force field along a path $L$ is the line integral
\begin{equation}
W=\int_L\mathbf{F}\cdot\mathrm{d}\boldsymbol{\ell}.
\end{equation}
For a gradient field, the chain rule gives
\begin{equation}
\int_A^B\nabla f\cdot\mathrm{d}\boldsymbol{\ell}=f(B)-f(A).
\end{equation}
Thus the line integral of a gradient depends only on its endpoints and vanishes around every closed path.

\begin{example}[Path independence]
Let $f=x^2y$, so that
\begin{equation}
\nabla f=2xy\,\mathbf{e}_x+x^2\,\mathbf{e}_y.
\end{equation}
For any path from the origin to $P=(x_0,y_0)$,
\begin{equation}
\int_O^P\nabla f\cdot\mathrm{d}\boldsymbol{\ell}
=f(P)-f(O)=x_0^2y_0.
\end{equation}
For the path that first follows the $x$ axis and then moves vertically, the first segment contributes zero and the second contributes
\begin{equation}
\int_0^{y_0}x_0^2\,\mathrm{d}y=x_0^2y_0.
\end{equation}
Any other path gives the same result because the field is a gradient.
\end{example}

\subsection{Flux and divergence}
The flux of a vector field $\mathbf{A}$ through a closed surface $S$ is
\begin{equation}
\Phi=\oint_S\mathbf{A}\cdot\mathrm{d}\mathbf{S},
\qquad
\mathrm{d}\mathbf{S}=\hat{\mathbf{n}}\,\mathrm{d}S,
\end{equation}
where $\hat{\mathbf{n}}$ is the outward unit normal. Positive net flux indicates a source inside the volume; negative net flux indicates a sink.

For a small rectangular volume, the net flux divided by the volume tends to
\begin{equation}
\nabla\cdot\mathbf{A}
=\frac{\partial A_x}{\partial x}
+\frac{\partial A_y}{\partial y}
+\frac{\partial A_z}{\partial z}.
\end{equation}
Thus divergence is the local rate of net outward flux per unit volume. The divergence theorem relates the local and global descriptions:
\begin{equation}
\oint_S\mathbf{A}\cdot\mathrm{d}\mathbf{S}
=\int_V(\nabla\cdot\mathbf{A})\,\mathrm{d}V.
\end{equation}

\begin{example}[Divergence theorem]
For $\mathbf{A}=x\mathbf{e}_x+y\mathbf{e}_y+z\mathbf{e}_z$,
\begin{equation}
\nabla\cdot\mathbf{A}=3.
\end{equation}
For a rectangular box with side lengths $a,b,c$ and one corner at the origin,
\begin{equation}
\int_V\nabla\cdot\mathbf{A}\,\mathrm{d}V=3abc.
\end{equation}
The face $x=a$ contributes $abc$ to the outward flux, the face $x=0$ contributes zero, and the $y$ and $z$ pairs contribute $abc$ each. Therefore the total surface flux is also $3abc$.
\end{example}

\subsection{Curl and Stokes' theorem}
The circulation of a vector field around a closed path $L$ is
\begin{equation}
C=\oint_L\mathbf{A}\cdot\mathrm{d}\boldsymbol{\ell}.
\end{equation}
The curl measures circulation per unit area. For a small loop in the $xy$ plane,
\begin{equation}
C\approx\left(\frac{\partial A_y}{\partial x}
-\frac{\partial A_x}{\partial y}\right)\Delta x\,\Delta y
=(\nabla\times\mathbf{A})_z\,\Delta S.
\end{equation}
More generally,
\begin{equation}
(\nabla\times\mathbf{A})\cdot\hat{\mathbf{n}}
=\lim_{\Delta S\to0}\frac{1}{\Delta S}
\oint_{\partial(\Delta S)}\mathbf{A}\cdot\mathrm{d}\boldsymbol{\ell}.
\end{equation}
The orientation of the loop and its normal are related by the right-hand rule.

Stokes' theorem follows by dividing a surface into small loops. Integrals on internal edges cancel because neighboring loops traverse each shared edge in opposite directions:
\begin{equation}
\oint_L\mathbf{A}\cdot\mathrm{d}\boldsymbol{\ell}
=\int_S(\nabla\times\mathbf{A})\cdot\mathrm{d}\mathbf{S}.
\end{equation}

\begin{example}[Stokes' theorem]
Let
\begin{equation}
\mathbf{A}=-y\,\mathbf{e}_x+x\,\mathbf{e}_y-z\,\mathbf{e}_z,
\end{equation}
and let $L$ be the circle of radius $R$ in the $xy$ plane. Along the circle,
\begin{equation}
\mathbf{A}\cdot\mathrm{d}\boldsymbol{\ell}=R^2\,\mathrm{d}\phi,
\qquad
\oint_L\mathbf{A}\cdot\mathrm{d}\boldsymbol{\ell}=2\pi R^2.
\end{equation}
Since $\nabla\times\mathbf{A}=2\mathbf{e}_z$, the disk integral is
\begin{equation}
\int_S(\nabla\times\mathbf{A})\cdot\mathrm{d}\mathbf{S}
=2\pi R^2,
\end{equation}
in agreement with the line integral. Any other oriented surface spanning the same boundary gives the same result.
\end{example}

The identities $\nabla\times(\nabla f)=\mathbf{0}$ and
$\nabla\cdot(\nabla\times\mathbf{A})=0$ are the local differential counterparts of these integral statements.